% !PRINTED 579
\weberpage{579}

%<<CONT
But this contradicts Theorem 5. of the preceding section, according to which $\tau$ can satisfy, with respect to a prime modulus $p$, no congruence of degree lower than $n$.

Accordingly it follows from (3) \gesperrt{that the fundamental number $\Delta$ of the field, apart from the sign, is the absolute norm of the function $F'(\tau)$}, and that therefore
\weq{5}{\pm \Delta = N_a[F'(\tau)]}
is to be set.

If now, instead of the $\omega_i$, another basis $\omega_i'$ of $\mathfrak{o}$ is taken as the foundation, then in place of $\tau$ there appears another form
\weq{6}{\tau' = \omega_1' t_1 + \omega_2' t_2 + \cdots + \omega_n' t_n,}
for which we may also set
\weq{7}{\tau' = \omega_1 t_1' + \omega_2 t_2' + \cdots + \omega_n t_n',}
and in these the $\omega_i'$ are connected with the $\omega_i$ by a linear substitution with the determinant $\pm 1$ (\S.\,145):
\weq{8}{(\omega_1', \omega_2', \ldots, \omega_n') = C(\omega_1, \omega_2, \ldots, \omega_n).}
If one carries out this substitution in (6), and arranges according to $\omega_1, \omega_2, \ldots, \omega_n$, there results
\weq{9}{(t_1', t_2', \ldots, t_n') = C_1(t_1, t_2, \ldots, t_n),}
where $C_1$ is the transposed substitution of $C$ (\S.\,37). If one now forms the function $F(t)$ for the function $\tau'$, there may result $F_1(t)$; let the derivative for $t=\tau'$ be $F_1'(\tau')$. If we now set
\[
F'(\tau) = \Psi(t_1, t_2, \ldots, t_n),
\]
then $\Psi$ is an integral function in $R$, and on account of (7) there results
\[
F_1'(\tau') = \Psi(t_1', t_2', \ldots, t_n').
\]

% [sic: The source prints $F'(\tau')$ as the first functional; the preceding definitions compare $F'(\tau)=\Psi(t_1,\ldots,t_n)$ with $F_1'(\tau')=\Psi(t_1',\ldots,t_n')$. The primed argument is retained as printed.]
Since the substitution (9) is reversible, it follows from this that the two functionals $F'(\tau')$ and $F_1'(\tau')$ are mutually divisible by one another (\S.\,143), and that they are consequently associated.

The function $F'(\tau)$, whose absolute norm is equal to the absolute value of the fundamental number, we therefore call the \gesperrt{fundamental functional}, and the ideal represented by $F'(\tau)$ the \gesperrt{fundamental ideal} of the field $\Omega$. Now let $\mathfrak{p}^e$ be the highest power of the prime ideal $\mathfrak{p}$ which divides $p$, and let $e \geqq 1$; then by 5. of the preceding section we can
\weq{10}{F(t) \equiv P(t)^e\, \Phi(t)\ (\mathrm{mod}\ p)}
set, where $\Phi(t)$ is an integral function in $R$ of such a nature that $\Phi(\tau)$ is not divisible by $\mathfrak{p}$. From (10), however,
%CONT>>
